\documentclass[11pt]{article}
\usepackage[T1]{fontenc}
\usepackage[a4paper,margin=25mm]{geometry}
\usepackage{amsmath,amssymb,amsthm}
\usepackage[hidelinks]{hyperref}
\newtheorem{theorem}{Theorem}
\title{A one-sided prime-interval test for all Dirichlet twists of a cubic Artin factor}
\author{Mathematical research note produced by GPT-6.2 Astra}
\date{22 November 2026}
\begin{document}
\maketitle
\begin{abstract}
For fixed $M$, consider primes $p\equiv1\bmod M$ in $(y,2y]$.
Count the roots of $X^3-269$ modulo each such prime, subtract one,
and weight by $\log p$. A lower bound of order $-y^\theta$ throughout
arbitrarily late intervals $[Y,Y^{1+\delta}]$ for every
$\theta>1/2$ forces the Riemann hypothesis for every Dirichlet twist
modulo $M$ of the cubic Artin factor. Conversely, RH for these twists
bounds the prime sum in absolute value by $O_M(\sqrt y\log^2(2y))$.
We give the reduction to Alp\"oge's general Landau-window theorem and
derive the converse from the individual Hecke explicit formula. The arithmetic lower bound
remains to be proved.
\end{abstract}

\section*{Statement}
Let $F=\mathbb Q(\sqrt[3]{269})$, $E=F(\zeta_3)$, $A=\mathbb Q(\zeta_M)$
and $D=\operatorname{Gal}(A/\mathbb Q)$. Write $\rho$ for the
two-dimensional standard representation of $\operatorname{Gal}(E/\mathbb Q)=S_3$.
For $p\nmid3\cdot269M$ put $t(p)=\#\{x\bmod p:x^3=269\}-1$ and define
\[
 B_M(y)=\sum_{\substack{p\le y,\ p\equiv1\bmod M\\p\nmid3\cdot269M}}
              t(p)\log p,\qquad T_M(y)=B_M(2y)-B_M(y).
\]
At $p\equiv2\bmod3$ the summand vanishes. At $p\equiv1\bmod3$ it has
coefficient $2$ or $-1$ according as $269$ is or is not a cube modulo $p$.

\begin{theorem}
Fix $M\ge1$ and $\delta>0$. Suppose that for each $1/2<\theta<1$
there are a constant $C_\theta$ and arbitrarily large $Y$ such that
\[
 T_M(y)\ge-C_\theta y^\theta \qquad(Y\le y\le Y^{1+\delta}).        \tag{1}
\]
Then every zero in $0<\Re s<1$ of every $L(s,\rho\otimes\chi)$,
$\chi\in\widehat D$, lies on $\Re s=1/2$. The same implication
holds with $B_M(y)$ in place of $T_M(y)$.

Conversely, RH for these twists implies
\[
   |B_M(y)|+|T_M(y)|\ll_M\sqrt y\log^2(2y)\qquad(y\ge2).
\]
\end{theorem}

\begin{proof}
Let $P(s)=\prod_{\chi\in\widehat D}L(s,\rho\otimes\chi)$,
$D_P=P'/P$ and $-D_P(s)=\sum b(n)n^{-s}$ for $\Re s>1$.
All factors are entire: $\rho$ is induced from a nontrivial cubic
character over $\mathbb Q(\zeta_3)$ and twisting preserves its
irreducibility. The induced twist is a nontrivial finite-order Hecke
character. The product $P=\zeta_{FA}/\zeta_A$ is nonzero at one by
the two simple-pole residues. It has real coefficients because the
sum of the underlying representations is $\rho\otimes{\rm reg}_D$.

At good primes $b(p)=\varphi(M)t(p)\log p$ if $p\equiv1\bmod M$
and is zero otherwise. At every prime power,
$|b(p^k)|\le2\varphi(M)\log p$ by the finite image on inertia
invariants. The finitely many excluded primes and the prime powers
with $k\ge2$ have total absolute weight
$O_M(\sqrt y\log^2(2y))$ up to $2y$. For example, summing
$\sum_{p\le(2y)^{1/k}}\log p$ over $2\le k\le\log_2(2y)$ gives this
bound directly.

Put $g(u)=-\sum_{e^u<n\le2e^u}b(n)$. Each summand is real and
measurable, and the sum is locally finite. The elementary bound
$\sum_{n\le2y}|b(n)|=O_M(y\log(2y))$ supplies
$|g(u)|=O_M(e^{2u})$ for $u\ge0$. For $\Re s>1$ absolute convergence
permits termwise integration. Writing $y=e^u$, each coefficient
$b(n)$ occurs over $n/2\le y<n$, whose left endpoint is at least
one since $b(1)=0$. The boundary $n=2$ starts exactly at $y=1$
and has the same integral on inclusion or exclusion of that point.
Consequently the Laplace transform required for $\Re s>2$ is
\[
\begin{aligned}
\Phi(s)
 &= -\sum_{n\ge2} b(n)\int_{n/2}^n y^{-s-1}\,dy\\
 &= -\frac{2^s-1}{s}\sum_{n\ge2} b(n)n^{-s}
  =\frac{2^s-1}{s}D_P(s).
\end{aligned}
\tag{2}
\]
The multiplier $(2^s-1)/s$ has no zero in $\Re s>0$. On that half-plane
$\Phi$ is holomorphic off the locally discrete set of zeros of $P$.
A zero $z$ of multiplicity $m$ produces a simple pole with nonzero
residue $m(2^z-1)/z$. If $z>0$ is real its residue is positive.
Thus $g,\Phi$ and this zero set satisfy every hypothesis of
\texttt{PoleData} in the cited formal statement. Complex residues
need only be nonzero. There is no assumed sum over zeros.

For $U=\log Y$, (1) and the omitted-coefficient estimate bound
$g(u)\le C'e^{\theta u}$ throughout $[U,(1+\delta)U]$, after increasing
the constant to cover the factor $\varphi(M)$ and the smaller error.
Choose the nonnegative, continuous majorant $R(u)=C'e^{\theta u}$.
Its integral with weight $e^{-\theta u}$ over the window is $C'\delta U$,
which is at most $e^{\gamma U}$ for every fixed $\gamma>0$ once $U$
is large enough. This is precisely \texttt{WindowHypU}. The analytic
theorem \texttt{Deduction.landau\_window} therefore forces
$\Re z\le\theta$ for every positive-real-part zero $z$.
Let $\theta\downarrow1/2$. Every twist is entire, so its zeros occur
in the product. The functional equation for the conjugate twist in the
same family reflects any zero left of $1/2$ into the excluded region.

For the cumulative sum use $g(u)=-\sum_{n\le e^u}b(n)$, whose Laplace
transform is $D_P(s)/s$. The residue is now $m/z$, again nonzero
everywhere in $\Re z>0$ and positive at positive real zeros. All other
checks are unchanged.

For the converse write $L_\chi=L(s,\rho\otimes\chi)$ and
$-L_\chi'/L_\chi=\sum c_\chi(n)n^{-s}$. This is the primitive Hecke
function over $K=\mathbb Q(\zeta_3)$ of the finite-order character
$\psi(\chi\circ N_{K/\mathbb Q})$, where conjugation inverts the cubic
character $\psi$ and fixes the norm pullback. Thus this character is
nontrivial, and the explicit formula has no pole term at one.
For $x\ge2$ its once-integrated version is
\[
 \sum_{n\le x}c_\chi(n)(x-n)
 =-\sum_{\lambda}\frac{x^{\lambda+1}}{\lambda(\lambda+1)}
       +O_\chi(x\log(2x)).
\]
The sum runs over the nontrivial zeros. Shifting the Mellin contour of
the logarithmic derivative with kernel $x^{s+1}/(s(s+1))$ gives the
identity. The trivial zeros and the kernel poles at zero and minus one
contribute the stated remainder. See Greni\'e--Molteni,
\href{https://sites.unimi.it/molteni/research/papers-pdf/44-molteni-An_explicit_Chebotarev_density_theorem_under_GRH.pdf}{\emph{An explicit Chebotarev density theorem under GRH}},
proof of Theorem~4.1, the individual formula for $I_\chi(x)$ preceding
(4.3), and Theorem~A.1 for the zero count.

Take the forward difference from $x$ to $x+h$, $h=\sqrt x$, and divide
by $h$. Under RH for this $L_\chi$, zeros of height at most $x/h$ contribute
$O_\chi(\sqrt x\log^2(2x))$ by the derivative bound and the unit-height
zero count $O_\chi(\log(2+|t|))$. For greater heights the two endpoints
give $O_\chi((x^{3/2}/h)\log(2x)/(x/h))$. The divided remainder costs
$O_\chi(\sqrt x\log(2x))$. Removing the last partial interval costs the
same, since $|c_\chi(n)|\le2\log n$ for $n\ge2$ by the degree-two
finite-image Euler factors. Hence
$\sum_{n\le x}c_\chi(n)=O_\chi(\sqrt x\log^2(2x))$.
Sum over the finitely many $\chi$ and remove higher powers and bad primes
as above. This gives the assertion for $B_M$, and subtraction gives $T_M$.
\end{proof}

\paragraph{The full Dedekind statement.}
Writing $\eta$ for the sign character of $S_3$, Artin factorization gives
\[
 \zeta_{EA}^{[E\cap A:\mathbb Q]}
    =\prod_{\chi} L(s,\chi)L(s,\eta\chi)L(s,\rho\otimes\chi)^2 .
\]
Full Dedekind RH for $EA$ therefore additionally requires RH for the
cyclotomic factors on the right. Those factors occur in the full zeta
function; the prime statistic in the theorem singles out the induced twists.

\paragraph{The analytic input.}
See the pinned formal source
\href{https://github.com/CocoSgt/riemann-hypothesis-lean/blob/7cbf7a0c32fca1badf05931879a37c3f8a2b03b4/src/Development/AutomorphicArithmetic/DedekindLFunctions/LandauZeroFreeCriterion.lean}
{for \texttt{PoleData} and \texttt{WindowHypU}}, and
\href{https://github.com/CocoSgt/riemann-hypothesis-lean/blob/7cbf7a0c32fca1badf05931879a37c3f8a2b03b4/src/Development/AutomorphicArithmetic/DedekindLFunctions/LandauTheorems/landauWindow.lean}
{for \texttt{Deduction.landau\_window}, lines 83--84}.
The argument uses this theorem as an input. Each transform in (2) was
computed from the stated finite prime sum with its boundary convention.
\end{document}
